% Part: many-valued-logic
% Chapter: syntax-and-semantics
% Section: connectives

\documentclass[../../../include/open-logic-section]{subfiles}

\begin{document}

\olfileid{mvl}{syn}{con}

\olsection{Basa lan Konektif}

Logika proposisional klasik, uga akeh logika liyane, nggunakake
himpunan \emph{tetapan proposisional} lan \emph{konektif}.
Contone, unsur primitif kang kita gunakake yaiku:
\begin{enumerate}
\tagitem{prvFalse}{Tetapan proposisional kanggo !!{falsity}~$\lfalse$.}{}
\tagitem{prvTrue}{Tetapan proposisional kanggo !!{truth}~$\ltrue$.}{}
\item Konektif logis:
  \startycommalist
  \iftag{prvNot}{\ycomma $\lnot$ (negasi)}{}%
  \iftag{prvAnd}{\ycomma $\land$ (konjungsi)}{}%
  \iftag{prvOr}{\ycomma $\lor$ (disjungsi)}{}%
  \iftag{prvIf}{\ycomma $\lif$ (!!{conditional})}{}%
  \iftag{prvIff}{\ycomma $\liff$ (!!{biconditional})}{}%
\end{enumerate}
\iftag{defNot,defOr,defAnd,defIf,defIff,defTrue,defFalse,defEx,defAll}{%
Saliyane konektif primitif ing ndhuwur, kita uga nggunakake simbol
kang ditetepake minangka cekakan, kayata
\startycommalist
  \iftag{defNot}{\ycomma $\lnot$ (negasi)}{}%
  \iftag{defAnd}{\ycomma $\land$ (konjungsi)}{}%
  \iftag{defOr}{\ycomma $\lor$ (disjungsi)}{}%
  \iftag{defIf}{\ycomma $\lif$ (!!{conditional})}{}%
  \iftag{defIff}{\ycomma $\liff$ (!!{biconditional})}{}%
  \iftag{defFalse}{\ycomma $\lfalse$ (!!{falsity})}{}%
  \iftag{defTrue}{\ycomma $\ltrue$ (!!{truth})}.}{}

Konektif kang padha uga digunakake ing logika mawa akeh nilai.
Nanging asring migunani yen ing siji logika ana pirang-pirang
versi, umpama, konjungsi. Supaya versi-versi iku bisa dibedakake,
simbolé kudu beda. Sawetara logika mawa akeh nilai uga ngemot
konektif kang ora nduweni padanan ing logika klasik. Mula kita
nganggo kerangka kang rada luwih umum tinimbang lumrahe.

\begin{defn}
  \emph{Basa proposisional} dumadi saka himpunan $\Lang L$
  kang isine \emph{konektif}. Saben konektif $\star$ nduweni
  \emph{aritas}; konektif kanthi aritas~$n$ diarani
  \emph{nduweni $n$ panggonan.}
  Konektif kanthi aritas~$0$ uga diarani \emph{tetapan};
  konektif kanthi aritas~$1$ diarani \emph{unar}, lan
  konektif kanthi aritas~$2$ diarani \emph{binar}.
\end{defn}

\begin{ex}
  Basa baku logika proposisional $\Lang L_0$ ngemot
  konektif iki (kanthi aritas saben-sabene):
  $\lfalse$~($0$),
  $\lnot$~($1$),
  $\land$~($2$),
  $\lor$~($2$),
  $\lif$~($2$). Akeh-akèhé logika kang kita rembug
  nggunakake basa iki. Sawetara logika, amarga tradhisi lan
  konvensi, nggunakake simbol liya kanggo sawetara konektif.
  Contone, ing logika produk, simbol konjungsi asring
  $\odot$ tinimbang~$\land$. Kadhang kala migunani kanggo
  nambah operator anyar, umpama operator katemtuan $\triangle$
  (nduweni $1$ panggonan).
\end{ex}

\end{document}
